% =====================================================================
\section{Why comparisons disagree}
\label{sec:intro}

Quantum computers were proposed to simulate quantum systems that classical computers cannot~\cite{feynman1982,lloyd1996}, and chemistry and materials science have long been regarded as their leading scientific applications~\cite{aspuru2005,reiher2017,bauer2020}. Classical computing has not stood still. Tensor-network methods reach strongly correlated systems once thought out of range~\cite{white1992,chan2011}. Neural networks represent many-body wavefunctions directly~\cite{carleo2017,carleo2019}. Learned models predict protein structures, density functionals and stable crystals at scale~\cite{jumper2021,luise2025,merchant2023}. By 2024 this progress had led some researchers to ask whether quantum computers will be needed for physics and chemistry simulation at all~\cite{gent2024}, and one quantum hardware group now advises choosing applications in which machine learning would struggle to generalize from a few quantum-computed solutions~\cite{babbush2025}.

The question this paper asks is narrower and, we think, more useful:
\begin{quote}
\emph{For a specified scientific answer, what evidence shows that one available or proposed computation is more useful than another?}
\end{quote}
Published comparisons often cannot answer it. A quantum advantage claim and its classical rebuttal usually appear separately, report different outputs at different accuracies, and state costs in units that do not convert: physical qubits, logical qubits, gate counts, core-hours, GPU-days. A fault-tolerant resource estimate is compared with a classical baseline chosen by the authors of the estimate. And in public discussion ``AI'' often names numerical methods that involve no learning. Four short examples show the kinds of mismatch that recur.

\begin{exbox}{Example 1. Reproducing a result is not always reproducing a computation}
A quantum sampling experiment is usually checked with a statistic such as a cross-entropy score. A classical calculation that matches the score has answered a validation question. It has not necessarily produced new samples from the full output distribution. Before asking which computer was faster, identify the output each side actually produced.
\end{exbox}
\begin{exbox}{Example 2. A model runtime is not a complete scientific answer}
A chemistry estimate may predict a few hours for a fault-tolerant circuit on a machine with about a million physical qubits. That time belongs to one prepared state, one circuit and one execution schedule. The useful answer may also need reliable state preparation, repeated executions and more than one electronic state. The classical result it is compared with may rest on an empirical extrapolation. Both are informative; together they are not yet a speed comparison at matched confidence.
\end{exbox}
\begin{exbox}{Example 3. A successful alternative workflow need not replace a solver}
An AI-assisted materials screen can identify promising compositions and guide successful experiments without solving the many-electron problem proposed for a quantum computer. Its proper comparison is the quality, cost and reliability of the resulting decision, including the data and validation it needs.
\end{exbox}
\begin{exbox}{Example 4. A speedup in one step is not a speedup of the job}
A quantum linear solver prepares a state proportional to the solution of $Ax=b$, and for some inputs its running time grows only logarithmically with the number of unknowns. That holds once the access to the matrix and right-hand side, the sparsity or encoding scale, the condition number, the precision and the wanted output are specified. The condition number measures how strongly an error in the input can be amplified in the solution; the cost grows at least linearly with that of the matrix as encoded, which preconditioning can reduce at a price. An engineering simulation must also load its data and read numbers out of the state. For the engineering cases reviewed, no practical advantage over a strong classical workflow is established, and the published estimates are far from practical (Section~\ref{sec:linear}).
\end{exbox}

\paragraph{What this paper contributes.}
\begin{enumerate}[leftmargin=*]
\item \textbf{An evidence map with an executable ledger.} We assemble \nEntries{} comparisons from 2017 to 2026 in which quantum hardware or a fault-tolerant resource estimate and a classical or AI method address the same or a closely related task. Each is recorded with the output it concerns, the kind of agreement reached, the method that reached it and the source version it was read from. The tables and scripts that generate every plotted number are released with the paper (Appendix~\ref{app:ledger}).
\item \textbf{Explanations and models.} We group the reasons classical methods succeed into four kinds (Section~\ref{sec:why}) and give explicit models for the quantities on which a crossover depends: complete cost on both sides, the success probability of phase estimation, the break-even point of a polynomial speedup, and the cost of input and output (Section~\ref{sec:cross}).
\item \textbf{Conditional assessments, positive and negative.} For \nClasses{} problem classes we state what is established, what is not, and which specific result would change the assessment (Section~\ref{sec:opps}, Appendix~\ref{app:forward}). The assessments include candidates for quantum value and one family, sparse linear solves for engineering simulation, where theory and published estimates point the other way.
\end{enumerate}

\paragraph{Scope.} This is a comparative survey. We run no new physical simulations or hardware benchmarks; reported values are reconciled, and scenario arithmetic is provided and labelled as such. The selection is curated, not exhaustive (Appendix~\ref{app:search}), so counts describe this evidence map and not the distribution of opportunities across science. No assessment is an impossibility result, and none carries a date.

\paragraph{Relation to earlier reviews.} LaRose gives a chronological account of advantage experiments and their refutations~\cite{larose2026}, and Hangleiter reviews the status of sampling claims~\cite{hangleiter2026}. Babbush et al.\ and Hoefler, H\"aner and Troyer analyze which speedups can pay off on error-corrected hardware~\cite{babbush2021,hoefler2023}. Dalzell et al.\ survey quantum algorithms and their end-to-end complexities~\cite{dalzell2023}, Aaronson sets out the input and output conditions that can erase linear-algebra speedups~\cite{aaronson2015}, and Lee et al.\ find no evidence of a generic exponential advantage in ground-state chemistry~\cite{lee2023}. Preskill, Bharti et al.\ and Daley et al.\ cover near-term devices and quantum simulation~\cite{preskill2018,bharti2022,daley2022}. We build on these. What we add is the pairing of each claim with the classical method that addressed it, at a stated level of task equivalence, together with resource accounting that keeps units and success semantics apart, and a record that can be corrected line by line.

\paragraph{How to read this paper.} Section~\ref{sec:compare} explains how we compare and ends with a one-table summary. Sections~\ref{sec:devices} and~\ref{sec:ft} review device experiments and fault-tolerant estimates. Section~\ref{sec:ai} describes what AI changes on the classical side. Section~\ref{sec:cross} gives the models, Section~\ref{sec:opps} the opportunities and the evidence each needs, and Section~\ref{sec:end} the limitations and conclusions. Technical terms are explained where they first matter and collected in Appendix~\ref{app:glossary}. Entries are cited by stable identifiers (E01--E\nEntries) that link the text to the appendix tables and the ledger. Every figure distinguishes its series by marker shape, line style, hatching or a printed label as well as by colour, and the captions identify series by those features and never by a colour name. The figures can therefore be read in greyscale, with any form of colour-vision deficiency, and in a viewer that renders colours differently.
